\section{Subexponential growth faster than every power}
\label{sec:superpolynomial}

The interpolation estimate completes the first subexponential branch.
We will find two large distinct clusters at distance whose logarithm
is negligible compared with the logarithm of their sizes. Interpolation
then places them across a single edge at a probability incompatible
with Proposition~\ref{prop:two-cluster}.

\begin{theorem}\label{thm:superpoly}
Let $G$ be an infinite connected locally finite quasi-transitive graph
with $\pc<1$. Suppose that
\[
 \log V_+(r)=o(r),
 \qquad \frac{\log V_-(r)}{\log r}\longrightarrow\infty.
\]
Then critical Bernoulli bond percolation on $G$ has almost surely no
infinite cluster.
\end{theorem}

\begin{proof}
Suppose instead that critical percolation has an infinite cluster with
positive probability. Let $\theta>0$ be the uniform lower bound in
\eqref{eq:uniform-jump}, and let $D$ be the maximum degree.
By \eqref{eq:positive-threshold}, $\pc>0$.
All constants below may depend on $G$ and $\theta$, but will be
independent of the scales and of $p\in[\pc/2,\pc)$.

\paragraph{A band of finite cluster sizes.}
Choose an arbitrarily large integer $L\ge2$. The event
$\{|C_v|\ge L\}$ is determined by a finite ball: if a cluster has
at least $L$ vertices, a rooted spanning exploration finds $L$ of them
within distance $L-1$ of $v$. Finite-event continuity and the finitely
many orbits therefore allow a single $p\in[\pc/2,\pc)$ for which
\[
 F_L(v):=\PP_p(|C_v|\ge L)\ge3\theta/4
 \qquad\text{for every }v.
\]
All clusters at this $p$ are finite almost surely. We seek a size
$s\ge L$ at which every vertex still has cluster-tail probability at
least $\theta/2$, while some vertex has a cluster in the band
$[s,s^2)$ with probability at least $(\log s)^{-2}$. The first condition
will force connections between nearby vertices; the second will supply
a finite cluster small enough to find another large cluster nearby.

Set $s_j=L^{2^j}$. The bands $[s_j,s_j^2)$ are disjoint, and the required
lower bounds $(\log s_j)^{-2}$ have a summable total. Let $k\ge1$ be
the first index for which
$F_{s_k}(z)<\theta/2$ for some orbit representative $z$.
Such an index exists because the finite-cluster tails decrease to zero.
For this $z$,
\[
 \sum_{j=0}^{k-1}\PP_p(s_j\le |C_z|<s_j^2)
 =F_L(z)-F_{s_k}(z)>\theta/4.
\]
On the other hand,
$\sum_{j\ge0}(\log s_j)^{-2}=4/(3(\log L)^2)<\theta/4$
when $L$ is sufficiently large. Hence there are $j<k$ and $s=s_j$
such that
\begin{equation}\label{eq:superpoly-band}
 f:=\PP_p(s\le |C_z|<s^2)\ge(\log s)^{-2},
 \qquad F_s(v)\ge\theta/2\quad\text{for every }v.
\end{equation}
In particular $s\to\infty$ as $L\to\infty$.

\paragraph{A nearby pair of large distinct clusters.}
Let $r$ be the least integer with $V_-(r)\ge s^4$.
Such an integer exists: an infinite connected locally finite graph has
vertices at every distance from each vertex, so $V_-(r)\ge r+1$.
The growth hypothesis implies
\begin{equation}\label{eq:superpoly-radius}
 \log r=o(\log s).
\end{equation}
Indeed, for every fixed $M>0$, eventually $V_-(t)\ge t^M$.
If $r-1$ is beyond this threshold, minimality gives
$(r-1)^M\le V_-(r-1)<s^4$, so $r<1+s^{4/M}$.
The alternative that $r-1$ remains below the threshold only improves
the bound. Letting $M$ be arbitrarily large proves
\eqref{eq:superpoly-radius}.

Condition on an exact incidence exposure of a band cluster
$K=C_z$, with $s\le|K|<s^2$. Its exterior has independent product law.
Resample all incidence edges of $K$ independently with parameter $p$,
leaving the exterior unchanged. For every fixed such exposure, the
resampled full configuration has law $\PP_p$, so the expected number
of vertices in $B(z,r)$ belonging to clusters of size at least $s$
in that configuration is at least
$(\theta/2)|B(z,r)|$.

In the exterior graph $G\setminus K$, at most
\[
 |K|+D|K|(s-1)\le s^2+Ds^3
\]
vertices lie in $K$ or in exterior clusters of size less than $s$
that are adjacent to $K$. Every vertex whose resampled cluster has
size at least $s$, but whose exterior cluster has size less than $s$,
lies in this set: otherwise its exterior cluster is unaffected by
resampling incidence edges of $K$. Consequently, uniformly in the
fixed band exposure, the expected number of vertices of $B(z,r)$
in exterior clusters of size at least $s$ is at least
\[
 (\theta/2)|B(z,r)|-s^2-Ds^3
 \ge(\theta/4)|B(z,r)|
\]
for sufficiently large $L$, because $|B(z,r)|\ge s^4$.
In the original configuration these exterior clusters are distinct
from $C_z=K$. Averaging over the band event and then over $B(z,r)$
therefore gives a deterministic $y\in B(z,r)$ such that
\[
 \PP_p(C_y,C_z\text{ distinct},\ |C_y|,|C_z|\ge s)
 \ge(\theta/4)f.
\]

Put $c_0=\theta^2/8$, and choose the smallest distance $d$ of any pair
$x,z'$ satisfying
\begin{equation}\label{eq:superpoly-separated-pair}
 \PP_p(C_x,C_{z'}\text{ distinct},\ |C_x|,|C_{z'}|\ge s)
 \ge c_0 f.
\end{equation}
The preceding pair shows that $d\le r$. Distinct clusters rule out
$d=0$. Proposition~\ref{prop:two-cluster} rules out $d=1$ for all
sufficiently large $L$, since its uniform bound $Cs^{-1/2}$ is smaller
than $c_0(\log s)^{-2}$.
For any pair $u,v$ at distance less than $d$, Harris's inequality and
\eqref{eq:superpoly-band} give probability at least $\theta^2/4$ that
both cluster sizes are at least $s$. By the minimality of $d$, the
probability that they are also distinct is less than $c_0f\le c_0$.
Thus
\begin{equation}\label{eq:superpoly-short-connections}
 \PP_p(u\leftrightarrow v)\ge\theta^2/8
 \qquad(d_G(u,v)<d).
\end{equation}

\paragraph{A finite graph to which interpolation applies.}
Fix a geodesic from $x$ to $z'$ of length $d$. We now choose a finite
induced ball $F$ sufficiently large; the parameter $p$, size $s$,
pair, and path remain fixed during this choice.
First include the balls of radius $s$ around $x,z'$. Every configuration
in the infinite-graph event \eqref{eq:superpoly-separated-pair} then
has two distinct endpoint clusters of size at least $s$ in $F$ as
well: the size witnesses lie in those balls, and restriction cannot
join distinct clusters. Therefore the finite endpoint-event probability
$\delta_F$ satisfies the exact lower bound
\begin{equation}\label{eq:superpoly-finite-endpoints}
 \delta_F\ge c_0f.
\end{equation}
No limiting equality or loss in this bound is needed.

For each proper pair on the geodesic, its path-index separation is
less than $d$, so \eqref{eq:superpoly-short-connections} applies.
Connections in increasing finite induced balls increase to the full
event, since every witnessing path is finite. There are only finitely
many pairs, so enlarge $F$ until their finite-graph connection
probabilities are all at least $\theta^2/16$.

Finally, all clusters at this fixed subcritical $p$ are finite almost
surely. Enlarge $F$ further until, for each edge $e$ of the fixed
geodesic, the probability that either full endpoint cluster is not
contained in $F$ is at most $s^{-1}$. If both clusters are contained,
their restrictions equal the original clusters. Hence
\begin{equation}\label{eq:superpoly-finite-edge-error}
 \PP_p^F(H_s(e))
 \le\PP_p(H_s(e))+s^{-1}
 \le Cs^{-1/2}+s^{-1}
\end{equation}
for every path edge. Enlarging $F$ preserves the earlier endpoint
and connection bounds.

Apply Proposition~\ref{prop:interpolation} in $F$ with
$c=\theta^2/16$. Its constants are uniform: $p\ge\pc/2>0$ and
$1-p\ge1-\pc>0$. By \eqref{eq:superpoly-band},
\eqref{eq:superpoly-radius}, and
\eqref{eq:superpoly-finite-endpoints}, some path edge satisfies
\[
 -\log\PP_p^F(H_s(e))
 \le C'\bigl(\log(1/(c_0f))+\log r+1\bigr)
 \le C'\bigl(2\log\log s+\log r+C''\bigr)
 =o(\log s).
\]
But \eqref{eq:superpoly-finite-edge-error} gives
\[
 -\log\PP_p^F(H_s(e))
 \ge\tfrac12\log s-\log(C+1)
\]
for $s\ge1$, a contradiction as $L\to\infty$.
This proves the theorem.
\end{proof}
